\documentclass[uplatex,a4paper,12pt]{jsarticle}
\usepackage{amsmath,amssymb}
\usepackage{enumitem}

\begin{document}

\title{ニューラルネットワーク基礎 確認テスト 解答}
\author{}
\date{}
\maketitle

\section*{配点}
全10問，各10点，合計100点．60点以上を合格とする．

\section*{解答}

\begin{enumerate}

\item
\[
\boxed{(b)}
\]
MLPは，多層パーセプトロンのことであり，パーセプトロンを多層化したニューラルネットワークである．

\item
\[
\boxed{(c)}
\]
XOR問題は，1本の直線では分類できない線形分離不可能な問題である．

\item
\[
\boxed{(b)}
\]
式
\[
y=g(\mathbf{w}^{\top}\mathbf{x}+b)
\]
における $g$ は活性化関数を表す．

\item
\[
\boxed{(a)}
\]
ReLU関数は，
\[
\mathrm{ReLU}(x)=\max(x,0)
\]
で定義される．

\item
\[
\boxed{\text{隠れ層}}
\]
または
\[
\boxed{\text{中間層}}
\]

\item
\[
\boxed{(b)}
\]
多値分類問題では，出力を各クラスの確率として扱うため，ソフトマックス関数がよく用いられる．

\item
\[
\boxed{(a)}
\]
学習率 $\eta$ は，勾配降下法におけるパラメータの更新量を調整する値である．

\item
\[
\boxed{\text{誤差逆伝播法}}
\]
または
\[
\boxed{\text{バックプロパゲーション}}
\]

\item
\[
\boxed{(c)}
\]
学習率は，学習前に設定するハイパーパラメータである．一方，重み $W$ やバイアス $b$ は学習によって調整されるパラメータである．

\item
\[
\boxed{(b)}
\]
MLPは，隠れ層と活性化関数を用いることで，非線形な関数を近似できる．

\end{enumerate}

\end{document}